\section{Related work}

Transport learning methods learn a mapping between two distributions, where the learned model can transform a sample from one distribution into a sample from the other one.
Typically, one of the distributions is easy to sample (such as a Gaussian) and the other one is the data distribution that one is interested in modeling.
The learned mapping can either be deterministic or stochastic.
A thorough overview of related areas can be found in \citet{yang2024diffusion}.

\paragraph{Deterministic transport.}
Deterministic transport methods implement a change of variable, either explicitly or approximately, that can be used to uniquely map a sample from one distribution to the other.
The normalizing flow family~\citep{rezende2015variational,dinh2017density,kingma2018glow} of methods construct an explicit invertible model that realizes this map either in one step or a finite number of discrete steps.
Neural ODEs~\citep{chen2018neural,grathwohl2018scalable} generalize from discrete steps to a continuous time mapping by inferring and learning the gradient field for all times.
However, Neural ODEs are difficult to train due to the need for simulating the ODEs as part of the training.
Rectified flows, flow matching, and iterative denoising methods~\citep{liu2022flow,lipman2022flow,tong2023improving,heitz2023iterative,delbracio2023inversion} either implicitly or explicitly specify a continuous mapping and learn a model to approximate the continuous time mapping.
Similarly, probability flow ODEs~\citep{song2020score} learned by diffusion models~\citep{sohl2015deep} approximate an implicitly defined continuous mapping.
Our work is useful for flexible sampling from such pre-trained continuous time deterministic Gaussian flows, or more generally where the score function for all the marginal distributions is either provided or can be deduced from the learned flow model.

\paragraph{Stochastic transport.}
Stochastic transport methods learn a stochastic mapping, where a sample from one distribution gets stochastically mapped to a sample from the other.
Gaussian diffusion models are a salient example of such discrete~\cite{sohl2015deep,ho2020denoising} or continuous time~\cite{song2020score,kingma2021variational} mappings where one of the distributions is constrained to be Gaussian.
Several generalizations have been proposed that extend from Gaussian to more general families of distributions~\cite{yoon2024score}.
The stochastic interpolants framework~\citep{albergo2022building,albergo2023stochastic,ma2024sit} further generalizes to a larger family of distributions by introducing a random latent variable allowing efficient estimation of the score function at all times.
Our work is directly applicable to models learned with such methods where the score function is accessible and can be used to construct and explore a large family of samplers.
The convergence rates of diffusion models have been studied in \cite{chen2023score,benton2023linear} with respect to the number of data samples and dimensionality.
However, since our method does not require retraining, it does not affect these properties of the original training algorithms.

\paragraph{Schr\"odinger bridge and optimal transport.}
These methods consider a more general problem of learning transport maps with additional constraints.
k-Rectified flows~\cite{liu2022flow} provide an iterative procedure for tackling deterministic optimal transports for a family of costs, while the more general Schr\"odinger bridge problem, viewed as an entropy regularized optimal transport, is an active area of research \cite{shi2024diffusion,liu2023i2sb}.
Our work is complementary to these methods as we focus on flexible sampling, given the access to the score function for the marginal distributions.
